当前师生对使用附录~\ref{app:config} 中的配对 completion 协议。分数以百分数而非小数表示；差异为同一权重步、同一 benchmark 上的 Block3 减去 Token。缺失 boxed 标记比例和引擎长度终止比例使用每题全部八条响应；Avg@8 与 Pass@8 定义见式~\ref{eq:metrics}。不同权重的比较描述每种方法的一条训练轨迹，不能解释为独立重复实验。
% Generated; do not edit. Requires booktabs; Chinese requires CJK support.
\begin{table*}[t]
\centering
\small
\caption{Qwen3-4B completion，seed21，COMPLETE。分数单位为百分比，NA表示缺失。}
\begin{tabular}{lllrrrr}
\toprule
基准 & 方法 & 步数 & Avg@8 & Pass@8 & 缺box & 截断 \\
\midrule
MATH500 & Block3 Mean & 200 & 75.40 & 90.60 & 0.90 & 0.90 \\
MATH500 & Block3 Mean & 150 & 78.03 & 91.20 & 0.83 & 0.83 \\
MATH500 & Block3 Mean & 100 & 75.52 & 90.80 & 1.25 & 1.18 \\
MATH500 & Block3 Mean & 50 & 78.53 & 91.80 & 0.85 & 0.75 \\
MATH500 & Token OPD & 200 & 75.92 & 91.00 & 0.85 & 0.80 \\
MATH500 & Token OPD & 150 & 76.17 & 91.40 & 0.95 & 0.97 \\
MATH500 & Token OPD & 100 & 76.92 & 90.40 & 1.07 & 0.95 \\
MATH500 & Token OPD & 50 & 78.65 & 92.20 & 0.85 & 0.80 \\
\midrule
AIME24 & Block3 Mean & 200 & 13.33 & 33.33 & 12.08 & 11.25 \\
AIME24 & Block3 Mean & 150 & 15.42 & 33.33 & 9.58 & 8.75 \\
AIME24 & Block3 Mean & 100 & 15.00 & 30.00 & 13.75 & 12.08 \\
AIME24 & Block3 Mean & 50 & 14.17 & 30.00 & 7.08 & 6.25 \\
AIME24 & Token OPD & 200 & 16.25 & 30.00 & 7.92 & 6.67 \\
AIME24 & Token OPD & 150 & 12.50 & 26.67 & 10.42 & 8.33 \\
AIME24 & Token OPD & 100 & 12.92 & 26.67 & 12.08 & 9.17 \\
AIME24 & Token OPD & 50 & 17.92 & 36.67 & 10.42 & 8.75 \\
\midrule
AIME25 & Block3 Mean & 200 & 13.75 & 30.00 & 4.17 & 3.33 \\
AIME25 & Block3 Mean & 150 & 11.67 & 26.67 & 5.42 & 5.42 \\
AIME25 & Block3 Mean & 100 & 12.08 & 33.33 & 7.50 & 6.25 \\
AIME25 & Block3 Mean & 50 & 11.67 & 23.33 & 7.50 & 5.42 \\
AIME25 & Token OPD & 200 & 13.33 & 30.00 & 3.75 & 3.33 \\
AIME25 & Token OPD & 150 & 15.83 & 40.00 & 5.42 & 6.25 \\
AIME25 & Token OPD & 100 & 10.83 & 23.33 & 5.83 & 5.42 \\
AIME25 & Token OPD & 50 & 12.92 & 23.33 & 7.08 & 4.17 \\
\midrule
AMC23 & Block3 Mean & 200 & 46.23 & 74.70 & 4.52 & 4.07 \\
AMC23 & Block3 Mean & 150 & 49.70 & 77.11 & 3.46 & 3.01 \\
AMC23 & Block3 Mean & 100 & 47.29 & 74.70 & 3.92 & 3.01 \\
AMC23 & Block3 Mean & 50 & 50.30 & 74.70 & 4.37 & 3.31 \\
AMC23 & Token OPD & 200 & 49.25 & 77.11 & 3.61 & 3.01 \\
AMC23 & Token OPD & 150 & 49.40 & 77.11 & 2.56 & 2.86 \\
AMC23 & Token OPD & 100 & 49.25 & 78.31 & 2.56 & 2.56 \\
AMC23 & Token OPD & 50 & 48.49 & 78.31 & 3.16 & 3.01 \\
\bottomrule
\end{tabular}
\par\smallskip\footnotesize 缺box仅指无字面标记 \texttt{\textbackslash boxed}；截断指引擎 length 终止。
\end{table*}


配对题目 bootstrap 区间在每个 benchmark 内独立计算，两组共同重采样相同题目身份，并条件化于每题已保留的八条响应及已训练权重。它不重采样训练 seed，不估计教师不确定性，也不提供跨多个权重或 benchmark 的同时覆盖保证。百分位方法使用 10,000 次重采样，seed 为 20260921。置信区间跨过零不证明两种方法等效。

\begin{figure}[t]
\centering
\includegraphics[width=\linewidth]{\qwenfigures/qwen4_pass.pdf}
\caption{当前配对实验四个权重的Pass@8。各任务纵轴范围独立，同一子图内两种方法共享范围。每个点使用全部题目、每题八条响应，并非选择部分输出。}
\end{figure}

\paragraph{事后重叠敏感性分析。}
剔除附录\ref{app:config}中的四道训练池字符串重叠题后，MATH500剩余496题。Step200的Avg@8差值为$-0.58$个百分点（$[-1.66,0.50]$），完整benchmark为$-0.52$。Step50/100/150的差值符号也未改变。这不能排除语义重叠或上游模型训练污染，不能估计污染的因果作用，也不构成新的留出benchmark。完整500题结果仍为主要报告。
% Generated post-hoc sensitivity only; requires booktabs and CJK for Chinese.
\begin{table*}[t]
\centering
\scriptsize
\setlength{\tabcolsep}{4pt}
\caption{事后重叠敏感性分析：MATH500的496题子集 (COMPLETE)。}
\label{tab:qwen4overlapsensitivity}
\begin{tabular}{lrrrrrr}
\toprule
 & \multicolumn{3}{c}{Avg@8} & \multicolumn{3}{c}{Pass@8} \\
\cmidrule(lr){2-4}\cmidrule(lr){5-7}
步数 & Block3 & Token & $\Delta$ & Block3 & Token & $\Delta$ \\
\midrule
200 & 75.73 & 76.31 & -0.58 [-1.66, 0.50] & 90.73 & 91.33 & -0.60 [-2.62, 1.21] \\
150 & 78.40 & 76.56 & +1.84 & 91.53 & 91.53 & +0.00 \\
100 & 75.83 & 77.27 & -1.44 & 90.93 & 90.52 & +0.40 \\
50 & 78.93 & 79.08 & -0.15 & 91.94 & 92.34 & -0.40 \\
\bottomrule
\end{tabular}
\par\smallskip\footnotesize 剔除审计中空白折叠加casefold匹配的4题，包含2题精确匹配。仅为事后敏感性分析，不是新基准或主结果；官方全500题分数不变。分数为百分比；$\Delta$为Block3减Token的百分点差；NA表示缺少配对题目分数。 仅Step200方括号为496题内配对问题bootstrap的逐点95\%区间；条件于一对训练运行，非多训练种子CI。
\end{table*}

